\chapter*{Conclusion}
\addcontentsline{toc}{chapter}{Conclusion}

Ce rapport a presente une chaine complete allant de la perception evenementielle
a l'integration navigation sur robot mobile. Le materiel (DAVIS346 + LIMO), les
fondements theoriques, puis les choix d'implementation ROS 2 ont ete explicites
de maniere coherente avec l'article de reference.

La contribution principale est l'adaptation d'une approche de detection d'objets
dynamiques basee sur compensation IMU, segmentation temporelle et clustering,
puis son raccordement a Nav2 via l'injection d'obstacles dynamiques dans les
costmaps.

Les suites naturelles du travail sont:
\begin{itemize}
\item renforcer la robustesse en scenes fortement dynamiques et peu texturees;
\item raffiner le suivi temporel des objets detectes;
\item consolider l'evaluation quantitative (latence, precision detection,
impact navigation) sur des campagnes de tests plus longues.
\end{itemize}
